% !TEX root = ../1Main.tex
% Hauptteil - Template/Beispielkapitel

\chapter{Hauptteil / Implementierung}
\label{chap:hauptteil}

% ===================================================================
% ANLEITUNG: Dies ist ein Beispielkapitel für den Hauptteil.
% Passen Sie Titel und Struktur an Ihre Arbeit an (z.B. "Konzept",
% "Implementierung", "Analyse", "Methodik" etc.)
% ===================================================================

In diesem Kapitel wird der Hauptteil der Arbeit beschrieben. Je nach Typ der Arbeit kann dies ein Konzept, eine Implementierung, eine empirische Untersuchung oder eine theoretische Analyse sein.

\section{Konzeptioneller Ansatz}
\label{sec:konzept}

Beschreiben Sie hier Ihren konzeptionellen Ansatz. Nutzen Sie Abbildungen zur Veranschaulichung:

% Beispiel für eine Abbildung (Bild muss vorhanden sein)
% \begin{figure}[htbp]
%     \centering
%     \includegraphics[width=0.7\textwidth]{beispiel-architektur.pdf}
%     \caption{Systemarchitektur}
%     \label{fig:architektur}
% \end{figure}

% Auf die Abbildung verweisen:
% \autoref{fig:architektur} zeigt die Gesamtarchitektur des Systems.

\subsection{Anforderungsanalyse}

Die wichtigsten Anforderungen an das System sind:

\begin{table}[htbp]
    \centering
    \caption{Funktionale Anforderungen}
    \label{tab:anforderungen}
    \begin{tabular}{lp{8cm}l}
        \toprule
        \textbf{ID} & \textbf{Beschreibung} & \textbf{Priorität} \\
        \midrule
        FA-01 & Das System muss \dots & Hoch \\
        FA-02 & Das System soll \dots & Mittel \\
        FA-03 & Das System kann \dots & Niedrig \\
        \bottomrule
    \end{tabular}
\end{table}

\section{Architektur und Design}
\label{sec:architektur}

Beschreiben Sie die Systemarchitektur und wichtige Designentscheidungen.

\subsection{Komponenten-Übersicht}

Das System besteht aus folgenden Hauptkomponenten:
\begin{description}
    \item[Komponente A] ist verantwortlich für \dots
    \item[Komponente B] übernimmt die Aufgabe \dots
    \item[Komponente C] stellt die Schnittstelle zu \dots bereit
\end{description}

\subsection{Technologie-Stack}

Folgende Technologien werden eingesetzt:
\begin{itemize}
    \item \textbf{Backend:} Python/Java/Node.js \dots
    \item \textbf{Frontend:} React/Vue/Angular \dots
    \item \textbf{Datenbank:} PostgreSQL/MongoDB \dots
    \item \textbf{Tools:} Git, Docker, CI/CD \dots
\end{itemize}

\section{Implementierung}
\label{sec:implementierung}

Details zur Implementierung werden hier beschrieben.

\subsection{Zentrale Algorithmen}

Der Hauptalgorithmus funktioniert wie folgt:
\begin{enumerate}
    \item Initialisierung der Datenstrukturen
    \item Verarbeitung der Eingabedaten
    \item Anwendung des Algorithmus
    \item Ausgabe der Ergebnisse
\end{enumerate}

% Code-Beispiele gehören oft in den Anhang, siehe Code.tex
% Verweise auf Code im Anhang: \autoref{lst:beispiel-code}

\subsection{Datenmodell}

Das Datenmodell umfasst folgende Entitäten \dots

\section{Besondere Herausforderungen}
\label{sec:herausforderungen}

Bei der Umsetzung traten folgende Herausforderungen auf:
\begin{itemize}
    \item \textbf{Herausforderung 1:} Beschreibung und Lösung
    \item \textbf{Herausforderung 2:} Beschreibung und Lösung
    \item \textbf{Herausforderung 3:} Beschreibung und Lösung
\end{itemize}

\section{Zusammenfassung}
\label{sec:hauptteil-zusammenfassung}

In diesem Kapitel wurden folgende Aspekte behandelt:
\begin{itemize}
    \item Konzeptioneller Ansatz und Anforderungen
    \item Architektur und Design
    \item Details zur Implementierung
    \item Besondere Herausforderungen und deren Lösungen
\end{itemize}

% ===================================================================
% HINWEIS: Passen Sie dieses Kapitel komplett an Ihre Arbeit an!
% ===================================================================
